% =========================================================================
% Chapter: Introduction
% Dissertation: A Comparative Study of Vision-Based Perception Paradigms
%               for Urban Waste and Infrastructure Monitoring
%               (MSc, University of Malta, 2026)
% Written against repository commit cee7dc6a (inspected 11 July 2026);
% revised against commit 2d70be27 (inspected 16 July 2026): attribute track
% now four backbone families / 16-variant size-adaptation matrix.
% Citation style: natbib author-year (\citet / \citep); keys resolve in
% Documentation/Background-LiteratureReview/background_literature_references.bib
% (extended with nso2026municipalwaste and eurostat2026municipal).
% NOTE: this chapter deliberately reports no experimental results; findings
% are presented in Chapter 5 and assessed in Chapter 6.
% =========================================================================

\chapter{Introduction}
\label{ch:introduction}

\section{Problem Definition}
\label{sec:problem-definition}

Municipalities are continuously responsible for the condition of public space: domestic waste must be collected reliably from the kerbside, streets must remain clean and passable, and traffic signs must stay visible, legible and structurally sound. Knowing \emph{where} intervention is needed at any given moment still depends substantially on manual inspection, periodic patrols, citizen reporting and fragmented organisational workflows, in which individual inspectors interpret visual conditions and relay them through separate administrative channels. Such arrangements are labour-intensive and geographically limited: each staff-hour of inspection covers a small area, coverage is intermittent rather than continuous, judgements differ between inspectors and are difficult to standardise, and the overall process is reactive---problems are often recorded only after they have persisted long enough to attract a complaint. Scaling this model across hundreds of streets and dozens of localities is impractical without some form of automation.

These difficulties are amplified in Malta. The islands are densely developed, with narrow roads, continuous building frontages and high visual clutter in which waste presentations, street furniture and signage compete for limited space. Kerbside waste is placed directly on pavements, where bags are frequently grouped, overlapping and partially occluded by parked vehicles; traffic signs are mounted variously on poles, walls and fa\c{c}ades, are often small or distant in street-level imagery, and are exposed to coastal weather that accelerates fading and corrosion of an already ageing signage stock. Illumination varies sharply between harsh sunlight, deep shadow and dusk, and monitoring imagery is realistically captured with smartphones, on foot or from moving vehicles rather than with calibrated survey equipment. Institutional responsibility is itself distributed---kerbside collection is organised through local and regional councils, public cleansing is provided by the Cleansing and Maintenance Division, and road infrastructure is shared between Transport Malta, Infrastructure Malta and the councils \citep{wastecollection2026web,cmd2026web,localcouncilsact1993}---so timely, georeferenced visual evidence would be valuable to several actors simultaneously.

This dissertation addresses two central application problems within this setting. The first is the detection and categorisation of kerbside domestic waste in street-level imagery, so that the presence and type of waste presentations can be recorded automatically. The second is the detection and richer understanding of traffic-sign and roadside-sign infrastructure, extending beyond category recognition to attributes relevant to maintenance, such as physical condition. The two problems are technically distinct and together span several task formulations: object detection predicts the location and semantic category of objects in a full image; attribute classification predicts characteristics of an already localised sign; prompt-based and open-vocabulary systems attempt to locate objects from textual or visual prompts without conventional task-specific supervised training \citep{kirillov2023segment}; and vision-language or reasoning-oriented systems may generate structured localisation outputs through multimodal interpretation rather than through a conventional detector head \citep{nvidia2025cosmosplatform}.

Finally, the problem is not reducible to a single headline accuracy value. A perception system that is to be deployed by a municipality must also be judged on its robustness in visually difficult scenes, its class-level and error behaviour (false positives that waste inspection effort, false negatives that leave problems unrecorded), its annotation and training requirements, its inference latency, throughput, model size and memory demand, its suitability for edge hardware, its sensitivity to prompt phrasing where prompts are used, the availability and calibration of output confidence \citep{guo2017calibration}, its reproducibility and maintainability, the integrity of the data on which it was trained and evaluated, and the privacy constraints attached to street imagery under the General Data Protection Regulation \citep{gdpr2016}. Comparing perception paradigms for municipal use therefore requires a study designed around these dimensions jointly, not around accuracy alone.

\section{Motivation}
\label{sec:motivation}

The societal case for better municipal monitoring is direct. Reliable waste collection and clean public spaces affect public health, amenity and tourism; safe and legible traffic infrastructure affects every road user, and a deteriorated sign degrades exactly the information drivers need at safety-critical locations. Timely identification of uncollected waste and damaged signs enables evidence-based maintenance prioritisation, directing limited operational resources where observed conditions warrant intervention rather than where fixed schedules place them. The scale of the underlying waste problem in Malta is documented: official statistics record 621\,kg of municipal waste generated per resident in 2024---574\,kg when the tourist population is included in the denominator---against a European Union average of 517\,kg \citep{nso2026municipalwaste,eurostat2026municipal}. Malta thus generates substantially more municipal waste per person than the EU average while operating one of the densest urban fabrics in Europe, which makes the efficiency of kerbside collection and street cleansing a standing operational concern.

The research motivation arises from the distance between this operational need and the coverage of existing work. Waste-recognition research has concentrated on isolated objects under controlled backgrounds, on material classification for recycling, and on facility-level sorting; other strands address aerial litter or environmental debris rather than kerbside domestic collection \citep{yang2016trashnet,proenca2020taco,bashkirova2022zerowaste,abdallah2020artificial}. None of these settings represents sealed, colour-coded refuse bags grouped on cluttered pavements, and because bag colours and presentation conventions are policy-specific, datasets built elsewhere cannot simply be reused. Traffic-sign research is mature in category recognition and detection \citep{stallkamp2012man,ertler2020mapillary}, but established datasets do not jointly represent the physical condition, mounting type, shape and viewing angle of each sign---precisely the attributes a maintenance-oriented inventory requires---and systems trained on international sign taxonomies cannot emit Malta-specific categories such as bilingual street nameplates or roadside convex mirrors. More generally, models trained on foreign benchmarks transfer imperfectly to new visual environments \citep{torralba2011unbiased}, which argues for locally collected, locally annotated evidence.

At the same time, the space of candidate approaches has widened faster than the evidence about their relative merits. Supervised detectors require domain-specific annotation and retraining, but produce compact, fast models. Promptable and open-vocabulary foundation models promise to reduce or remove annotation requirements \citep{kirillov2023segment,meta2025sam3}, yet introduce prompt sensitivity, inconsistent localisation behaviour and, for generative variants, missing confidence calibration; large foundation models may also be too expensive for edge deployment \citep{chen2019deep}. Self-supervised representations may reduce the need for full end-to-end fine-tuning \citep{simeoni2025dinov3,hu2022lora}, but their suitability for locally specific sign attributes is untested. Published results for these families are produced on different datasets, hardware, image resolutions, training strategies and model scales, so their benchmark values cannot be compared directly. What the literature lacks---and what a municipality would need before committing to any approach---is a controlled, Malta-specific, deployment-aware comparison in which competing paradigms are evaluated on the same data, the same ground truth and the same operational criteria. Establishing that comparison is the purpose of this dissertation; the full critical treatment of prior work is given in Chapter~\ref{ch:literature-review}.

\section{Proposed Solution}
\label{sec:proposed-solution}

The dissertation constructs and evaluates a two-domain, three-paradigm comparative benchmark grounded in purpose-built Maltese datasets.

The first dataset, the \emph{Maltese Domestic Waste Dataset} (MDWD), is a street-level object-detection dataset representing kerbside domestic waste within Malta's municipal collection context. Its current working release comprises 3{,}598 unique source photographs annotated with five operational waste categories---Mixed Waste, Organic Waste, Recyclable Material, Orange CMD and Other Waste---that correspond to visually distinguishable waste presentations aligned with Malta's collection streams (the augmented training exports derived from these source images are larger and are described in Chapter~\ref{ch:methodology}; they are not unique captures). The second dataset, the \emph{Maltese Traffic Sign Dataset} (MTSD), is a street-level dataset of traffic-sign and roadside-sign infrastructure collected in eleven groups totalling 7{,}492 raw images. At the time of writing, six groups have passed annotation quality assurance, yielding 3{,}808 QA-verified images containing 10{,}263 annotated sign instances over twelve object classes spanning safety-relevant, regulatory, navigation-related and municipal signage, including explicit categories for rear-facing and otherwise unidentifiable sign infrastructure. Every annotated instance additionally carries four per-object attributes: viewing angle (front, side, back), mounting type (pole- or wall-mounted), physical condition (good, weathered, heavily damaged) and geometric shape (circular, quadrangle, triangular, octagonal, pentagon). Both datasets are characterised through systematic exploratory analysis, and their reliability is treated as part of the research problem: annotation quality auditing, dataset integrity checking and split-leakage analysis are built into the study rather than assumed away. Full taxonomy definitions appear in Chapter~\ref{ch:background}.

Three perception paradigms are compared on these datasets. The supervised track trains and evaluates object detectors from three recent YOLO generations---YOLO11, YOLO12 and YOLO26---and the RF-DETR detection transformer, representing contemporary one-stage convolutional, attention-oriented and transformer-based detection design across the model capacities each family supports. The prompt-based track evaluates foundation models that localise objects from natural-language prompts without task-specific training: SAM~3 and SAM~3.1, which perform promptable concept segmentation and return masks, boxes and confidence scores \citep{meta2025sam3,meta2026sam31}; LocateAnything, a dedicated visual-grounding vision-language model that decodes bounding boxes in parallel \citep{nvidia2026locateanything}; and NVIDIA Cosmos Reason2 in its 2B and 8B variants, with the 32B variant as an optional computationally heavy configuration---reasoning-oriented vision-language models from a world-foundation-model ecosystem that emit structured box coordinates through multimodal interpretation \citep{nvidia2025cosmosreason2,nvidia2025cosmosplatform}. These systems are deliberately not treated as interchangeable: they differ in output type, in whether confidence scores exist, and in mechanism, and the study assesses them as zero-shot alternatives or complements to supervised detection rather than as drop-in equivalents. The third track addresses MTSD attribute classification through a shared-backbone, multi-head architecture in which one visual encoder feeds four independent attribute heads. Sixteen variants form a size/adaptation matrix over four pretrained representation families at base and large scales: DINOv3, a general-purpose self-supervised image transformer, frozen-probed and LoRA-adapted \citep{simeoni2025dinov3,hu2022lora}; V-JEPA~2.1, a video-pretrained joint-embedding predictive architecture, likewise frozen-probed and LoRA-adapted \citep{murlabadia2026vjepa21}; ConvNeXt, a modern supervised convolutional network, frozen-probed and fully fine-tuned \citep{liu2022convnet}; and LingBot-Vision, a boundary-centric self-supervised transformer pretrained for dense spatial perception, frozen-probed and LoRA-adapted \citep{fu2026lingbot}. An earlier six-variant round over three families at a single size each is retained as a historical record. This is an adaptation-strategy and model-size comparison across different pretrained representation objectives---semantic self-distillation, temporal predictive embedding, supervised classification and boundary-centric dense pretraining---rather than a strictly controlled universal backbone ranking.

The evaluation is designed to answer deployment questions, not only accuracy questions. Detection is assessed with precision, recall, F1, average precision and mean average precision at standard intersection-over-union thresholds, together with per-class behaviour; attribute classification is assessed with accuracy, macro- and per-class F1 and confusion structure, reflecting pronounced class imbalance; robustness is examined across visually difficult slices, supported by qualitative failure cases and statistical uncertainty estimates where applicable. Alongside these, the study records parameter counts, checkpoint sizes, inference latency and throughput, annotation and training requirements, prompt sensitivity, edge-deployment suitability, maintainability and dataset integrity. Because the paradigms do not produce equivalent outputs---some prompted and reasoning-oriented systems emit no calibrated per-box confidence, so ranking-based metrics such as average precision cannot be computed for them meaningfully---the methodology employs complementary quantitative and qualitative measures appropriate to each output type, with all metric definitions, thresholds and protocols specified in Chapter~\ref{ch:methodology} and all findings reported in Chapter~\ref{ch:evaluation}.

The principal contributions of the dissertation follow from this design. First, it develops and systematically characterises two Malta-specific street-level datasets---MDWD for kerbside domestic waste and MTSD for traffic-sign infrastructure---with operationally grounded object taxonomies and, for MTSD, a per-sign attribute schema (condition, mounting, viewing angle, shape) that established traffic-sign datasets do not jointly provide. Second, it constructs a reproducible comparative benchmark of contemporary supervised detectors spanning three YOLO generations and a detection transformer on the waste-detection task. Third, it integrates promptable segmentation, grounded vision-language localisation and world-foundation-model reasoning into a single fixed-protocol evaluation against municipal ground truth, making structurally dissimilar systems comparable to supervised references on common data. Fourth, it contributes a shared-backbone multi-head attribute-classification framework comparing frozen probing, LoRA and full fine-tuning across four pretrained representation families and two model sizes. Fifth, it delivers a deployment-oriented evaluation framework---covering robustness, annotation effort, efficiency, edge suitability and dataset reliability---together with the audited datasets, protocols and tooling that allow the benchmark to be extended or reproduced.

\section{Aims and Objectives}
\label{sec:aims-objectives}

The overall aim of this dissertation is to systematically evaluate the suitability of contemporary vision-based perception paradigms for automated kerbside-waste and traffic-sign infrastructure monitoring in Malta, considering predictive performance, robustness, computational efficiency and operational practicality. This will be achieved through the following objectives:

\begin{enumerate}
  \item To construct, quality-assure and statistically characterise two Malta-specific street-level datasets---MDWD for kerbside domestic-waste detection and MTSD for traffic-sign infrastructure with per-sign condition, mounting, viewing-angle and shape attributes---including annotation auditing, dataset-integrity checking and split-leakage analysis.
  \item To train and quantitatively compare supervised object detectors from the YOLO11, YOLO12 and YOLO26 generations and the RF-DETR detection transformer on MDWD across the supported model capacities, using mean average precision, per-class and error-behaviour measures under a common training and evaluation protocol.
  \item To evaluate the zero-shot localisation capability and operational constraints of promptable and reasoning-oriented foundation models---SAM~3, SAM~3.1, LocateAnything and Cosmos Reason2---against fixed municipal ground truth under a documented prompt protocol, quantifying detection quality, prompt sensitivity and comparability limits relative to the supervised reference systems.
  \item To compare frozen linear probing, LoRA and full fine-tuning across DINOv3, V-JEPA~2.1 and ConvNeXt backbones for shared-backbone multi-head traffic-sign attribute classification, and to assess cross-paradigm deployment trade-offs in annotation requirements, model size, latency, throughput and edge suitability.
\end{enumerate}

Each objective is evaluated later in the dissertation. Objective~1 is assessed through the exploratory analyses, annotation-quality audits, integrity reports and leakage-sensitivity analysis reported for both datasets. Objective~2 is assessed through the MDWD detection benchmark, using mean average precision at standard thresholds, per-class metrics and error analysis. Objective~3 is assessed through the fixed-protocol prompt-based evaluation on the same ground truth, using operating-point precision, recall and F1 together with prompt-sensitivity and failure-case analysis, with average precision reported only where calibrated confidences exist. Objective~4 is assessed through the six-variant attribute-classification comparison using macro- and per-class F1 and confusion analysis, and through the deployment measurements of model size, latency and throughput. The corresponding procedures are specified in Chapter~\ref{ch:methodology} and the outcomes are reported in Chapter~\ref{ch:evaluation}; whether, and to what extent, each objective was achieved is assessed in Chapter~\ref{ch:conclusion}.

The scope of what this study can establish should be stated plainly at the outset. The dissertation is designed to determine how the selected paradigms compare on two specific Maltese municipal tasks, on datasets whose size, geographic coverage and annotation scope are finite and documented; it does not attempt universal claims about model families, and differences in pretraining data, model scale and output type mean that some comparisons are necessarily conditional rather than absolute. These boundaries, and their consequences for interpretation, are made explicit where the evidence is presented.

\section{Document Structure}
\label{sec:document-structure}

This dissertation progresses from the underlying municipal-monitoring and computer-vision context to the design, implementation and empirical evaluation of the proposed comparative benchmark, and is organised as follows. Chapter~\ref{ch:background} establishes the conceptual and domain context: municipal monitoring and the Maltese institutional setting, the MDWD and MTSD domains and their object and attribute taxonomies, the computer-vision task formulations used throughout the work, and family-level introductions to object detectors, foundation models, representation learning, multi-head classification, edge computing and evaluation terminology. Chapter~\ref{ch:literature-review} critically reviews prior work---municipal monitoring applications, waste and traffic-sign datasets, the evolution of supervised and end-to-end detection, promptable and open-vocabulary localisation, reasoning-oriented foundation models, the DINOv3, V-JEPA~2.1 and ConvNeXt representation families, adaptation strategies and multi-task attribute learning---closing with the research gaps this dissertation addresses. Chapter~\ref{ch:methodology} details how the study was conducted: dataset construction, annotation and quality assurance, detector training, the prompt-based evaluation protocol, the attribute-classification pipeline, experimental controls, metric definitions, robustness and uncertainty analyses, inference benchmarking and reproducibility measures. Chapter~\ref{ch:evaluation} presents and critically analyses the results of all three experimental tracks, including model-capacity and class-level comparisons, robustness and failure-case analysis, efficiency measurements and the operational trade-offs between paradigms. Finally, Chapter~\ref{ch:conclusion} revisits the aim and objectives, synthesises the findings, states the contributions, acknowledges the limitations of the study, and outlines future work, including broader geographical coverage, additional annotated MTSD collection groups, refined physical-condition assessment and progression towards practical municipal deployment. The next chapter begins this progression by establishing the municipal context and technical foundations on which the remainder of the dissertation builds.
